\begin{center}\begin{tikzpicture}[font=\small,>=stealth,line join=round]
\draw[gray!40,rounded corners=3pt] (0,0) rectangle (14.2,10.0);
\node at (7.1,9.4) {补充边界：两点都在轴上时，最短位置不再唯一};

\begin{scope}[shift={(2,2)},x=1.0cm,y=1.0cm]
\draw[->,gray] (-1,0)--(5.6,0) node[right]{$x$};\draw[->,gray] (0,0)--(0,5.8) node[above]{$y$};
\draw[blue,very thick] (-1,6)--(1,2)--(3,2)--(5,6);
\draw[orange!85!black,very thick] (1,2)--(3,2);\fill[orange!85!black] (1,2) circle(2pt) (3,2) circle(2pt);
\draw[dotted,gray] (1,0)--(1,2) (3,0)--(3,2);\node[below] at (1,0){$1$};\node[below] at (3,0){$3$};
\end{scope}
\node[align=left,anchor=west] at (8.3,5){$f(x)=|x-1|+|x-3|$\\[12pt]$a=b=0$\\[12pt]所有 $x\in[1,3]$\\[9pt]都取得最小值 $2$\\[12pt]含 $a+b$ 分母的公式失效};

\node[text width=12.8cm,align=center] at (7.1,.65) {橙色整段都最短；此时不能使用唯一取等位置公式。};
\end{tikzpicture}\end{center}
